\section{Interpretation and product implications}
The study yields a conditional engineering result about evidence provision. Catalog-conditioned selection makes complete positive support available under the tested small input budget, and the same module composes successfully with the general-purpose GBrain retrieval component. Larger contexts close the observed support gap. Separately, many shared-reader failures remain, including recommendations that violate explicit policy or resource constraints. The original primary accuracy interval includes zero. These are two distinct findings: a more sufficient context is not a reliable or commercially useful recommendation by itself.

The central distinction is between storing evidence, selecting a sufficient set, and deciding correctly from it. Relevance ranking alone performs poorly on these deliberately repetitive distractors at a small allowance. Generic relation diversity closes the support gap when more context is available. Required-slot coverage is a plausible selection mechanism for an explicit catalog, but its algorithm is not uniquely marketing-specific. These results do not establish a graph advantage, a need for a new database, or superiority over full deployed memory systems.

For a marketing-memory implementation, the most direct implication is to preserve approved strategy and provenance, resolve explicit replacements, and check formal prerequisites outside the generative reader. The proposed action should retain its policy versions, evidence bindings, objective and resource allowance. A validator should distinguish a supported proposal, a violated constraint and missing evidence before any real execution. The SQL result and the reader's negative-case errors support this division of responsibility on the tested clean-record task. They do not establish how much extraction or policy-authoring work is needed to make it reliable in practice.

The next research boundary is therefore evidence quality and decision usefulness under less favorable inputs. Independently authored histories should vary catalog alternatives more substantially; raw-document ingestion should charge extraction and maintenance cost; and stronger readers or different interfaces should be compared on a new frozen split. Blinded GTM assessment could then measure whether eligible proposals are useful, specific and worth a human's time. Campaign outcomes would require a separate prospective evaluation. Those are unmeasured questions, not implied benefits of the present scores.

Full context remains a serious reference at this record size. Its slightly higher primary point estimate comes with more input and generation time, but the paired interval is too broad to claim equivalence or a reliable quality ordering. A product should consequently choose between full context and selection using its actual context sizes, failure costs and reader, rather than treating this experiment as a universal architecture prescription.

